\documentclass[10pt,a4paper]{amsart}
\usepackage[margin=19mm]{geometry}
\usepackage{amsmath,amssymb,mathtools}
\usepackage[scr=rsfs,cal=euler]{mathalfa}
\usepackage{xfrac}
\usepackage[unicode,hidelinks,bookmarks=false]{hyperref}
\newcommand{\ie}{i.e.\ }
\newcommand{\cf}{cf.\ }
\newcommand{\resp}{respectively\ }
\newcommand{\Subsection}{\subsection}
\providecommand{\nobreakdash}{\nobreak\hbox{-}}
\setlength{\emergencystretch}{2em}
\multlinegap0pt
\usepackage{bm}
\usepackage{bbm}
\usepackage{dsfont}
\usepackage{xspace}
\usepackage{cancel}
\usepackage{mathtools}
\usepackage{cleveref}
\usepackage{amsthm}
\makeatletter
\@ifundefined{lemma}{\newtheorem{lemma}{Lemma}}{}
\@ifundefined{theorem}{\newtheorem{theorem}{Theorem}}{}
\makeatother
\title[Rates of convergence in long time asymptotics of an alignmen]{Rates of convergence in long time asymptotics of an alignment model with symmetry breaking}
\author{Alexandre Surin}
\date{Source: arXiv:2509.06765v3 (CC BY 4.0)}
\begin{document}
\maketitle

\noindent\emph{Source and licence.} Rates of convergence in long time asymptotics of an alignment model with symmetry breaking. Version arXiv:2509.06765v3 (CC BY 4.0), distributed under the Creative Commons Attribution 4.0 International licence, as recorded in the arXiv licence metadata for this exact version and shown on the abstract page https://arxiv.org/abs/2509.06765v3. Full rights notice: licenses/arxiv-2509.06765-CC-BY-4.0.txt.\par\smallskip

\noindent\textit{Editorial scope.} Counted unit: the Lemma whose source label is \texttt{deriveps} in the article (source0.tex lines 1036--1047, source label \texttt{deriveps}), delivered together with the article's own complete original proof of it (source0.tex lines 1048--1100, one \texttt{proof} environment of 53 lines). No mathematics is written, repaired or re-derived: statement and proof are the article's own text line for line. Required context retained uncounted: PDE (lines 141--144); thm2 (lines 198--206); thm3 (lines 211--217); ustar (lines 750--753); definitionQ1 (lines 804--807); gstar (lines 979--982); evoldegstar (lines 989--992); nalblau (lines 1007--1010). Every definition, display and label of those lines is retained unchanged. Omitted from this package and not redistributed: 1--140, 145--197, 207--210, 218--749, 754--803, 808--978, 983--988, 993--1006, 1011--1035, 1101--1347. Nothing inside a retained excerpt is omitted or altered: every retained body is delivered exactly as the article's lines are written, so no omission receipt of any kind accompanies this package. Citations: the retained excerpts carry no citation command at all, so no bibliography entry of the article is transcribed and no reference is redistributed. Reference adaptation: none was needed and none was made; every reference inside the delivered unit resolves inside it, so no unresolved reference or citation is produced. Printed slips kept literal: no printed word, symbol, hypothesis or number of the counted statement or of its proof was altered. Layout and class: the article's own class is \texttt{amsart} with packages including inputenc, aeguill, babel, geometry, amsmath, amssymb, amsfonts, amsthm, amsxtra, mathrsfs, bbold, mathtools, dsfont, graphicx; the wrapper is the lane's pinned \texttt{amsart} editorial wrapper at 10pt on A4, which restates the theorem-like environments the retained text needs and adds the article's own macro definitions that the retained text needs (none), with extra wrapper packages (bm, bbm, dsfont, xspace, cancel, mathtools, cleveref). The delivered numbering is the unit's own. Assets: no retained span contains a figure environment, a graphics-inclusion command, a PDF-page inclusion, a table or a TikZ picture; no image file is copied into this package, the package declares no asset dependency and no image file is redistributed with it. Licence: version arXiv:2509.06765v3 of this article is distributed under the Creative Commons Attribution 4.0 International licence, as recorded in the arXiv licence metadata for this exact version and shown on the abstract page https://arxiv.org/abs/2509.06765v3. A full rights notice is delivered at licenses/arxiv-2509.06765-CC-BY-4.0.txt. Characters: every retained excerpt is pure ASCII, so the delivered engine, XeLaTeX with its default Latin Modern text font, reports no missing character.
\begin{equation}
\label{PDE}
\frac{\partial f}{\partial t} = D \, \Delta f + \nabla \cdot \big[ \left(\nabla \psi_{\alpha}(v)+v - \mathbf{u}_{f}   \right) f\big]\quad \text{with} \quad \psi_{\alpha}(v) = \frac{\alpha}{4} \,|v|^4 - \frac{\alpha}{2}\,|v|^2.
\end{equation}

\begin{theorem}{}
\label{thm2}
Let $D<D_*$. Let $f$ be a nonnegative solution to \Cref{PDE} with  initial \\ \mbox{datum ~$f_{\mathrm{\mathrm{in}} } \in L^1_+(\mathbb{R}^d)$} such that ~$\mathcal{F}[f_{\mathrm{\mathrm{in}}}]<\mathcal{F}[G_{0}]$ and ~$\| f_{\mathrm{in}}\|_{L^1(\mathbb{R}^d)}=1$. Then
   the average speed $\mathbf{u}_{f}$ converges as $t \rightarrow\infty$ to a unique limit $\mathbf{u}_{\infty} \in \mathcal{S}$ and for any $\beta>0$,
  \begin{align*}
\lim_{t\rightarrow \infty} t^{\beta}\left(\mathcal{H}\big[f(t,\cdot) \mid G_{\mathbf{u}_{\infty}}\big] + |\mathbf{u}_{f}(t) - \mathbf{u}_{\infty}|^2\right) = 0
\end{align*}

\end{theorem}

\begin{theorem}{}
\label{thm3}
There is a positive constant $\beta$ depending only on $\alpha$ and $D$ such that under the assumptions of \Cref{thm2}, if $f_{\mathrm{in}} \in \bigcap_{u \in \mathcal{S}} L^2(\mathbb{R}^d,G_\mathbf{u}^{-1}dv)$, then
     \begin{align*}
         \mathcal{H}\big[f(t,.)|G_{\mathbf{u}_{\infty}}\big] +|\mathbf{u}_{f}-\mathbf{u}_{\infty}|^2 =O(e^{-\beta t}) \text{ as } t\rightarrow +\infty.
     \end{align*}
\end{theorem}

\begin{align}
\label{ustar}
    \mathbf{u}_{*}(t)= r(D) \frac{\mathbf{u}_{f}(t)}{|\mathbf{u}_{f}(t)|}. 
\end{align}

 \begin{align}
 \label{definitionQ1}
    \langle g_1,g_2\rangle_{\mathbf{u}} = D \int_{\mathbb{R}^d}g_1g_2 G_{\mathbf{u}_{}}dv -D^2 v_{g_1}v_{g_2}.
    \end{align}

\begin{align}
\label{gstar}
    f=G_{*}(1+g),
\end{align}

     \begin{align}
     \label{evoldegstar}
         \partial_tg=\mathcal{L}_{*}g+\mathcal{R}_{*}g-\frac{v\cdot \mathbf{u}_*'}{D}(1+g)
     \end{align}

    \begin{align}
     \label{nalblau}
         \nabla_{\mathbf{u}}G_{\mathbf{u}}(v)=\frac{v-\mathbf{u}}{D}G_{\mathbf{u}}(v)
     \end{align}

 \begin{lemma}
\label{deriveps}
 Under the assumptions of \Cref{thm3}, if $f$ solves \eqref{PDE} and $g$ and $\mathbf{u}_*$ are defined by \eqref{gstar} and \eqref{ustar}, then 
    \begin{align}
        \frac{1}{2}\frac{d}{dt}\langle g,g\rangle _{\mathbf{u}_*}= -Q_{2,\mathbf{u}_*}[g]+ R[g]
    \end{align}
where 
\begin{align}
\label{expressiondeR}
    R[g]&=D^2 \mathbf{v}_g  \int_{\mathbb{R}^d}g(\nabla g- \mathbf{v}_g)G_{\mathbf{u}_*}\, dv -\frac{\mathbf{u}_*'}{2}\cdot \int_{\mathbb{R}^d}g^2vG_{\mathbf{u}_*}dv.
\end{align}
\end{lemma}

\begin{proof}
We have:
\begin{equation*}
\begin{split}
    \frac{1}{2}\frac{d}{dt}\langle g, g \rangle_{\mathbf{u}_*}
    &= \langle g, \partial_t g \rangle_{\mathbf{u}_*}
    + \frac{1}{2} \nabla_\mathbf{u} \langle g, g \rangle_\mathbf{u} \cdot \mathbf{u}_*'.
\end{split}
\end{equation*}
Using \eqref{evoldegstar}, we obtain,
\begin{equation}
\label{bigequa}
\begin{split}
    \frac{1}{2} \frac{d}{dt} \langle g, g \rangle_{\mathbf{u}_*}
    &= -Q_{2,\mathbf{u}_*}[g] 
    + D^2\, \mathbf{v}_g \cdot \int_{\mathbb{R}^d} g(\nabla g - \mathbf{v}_g) G_*\,dv \\
    &\quad - \frac{1}{D} \langle v \cdot \mathbf{u}_*'(1+g), g \rangle_{\mathbf{u}_*}
    + \frac{1}{2} \nabla_\mathbf{u} \langle g, g \rangle_\mathbf{u} \cdot \mathbf{u}_*'
\end{split}
\end{equation}
Let us start with the computation of $\langle v\cdot \mathbf{u}_*'(1+g),g\rangle_{\mathbf{u}_*}$.
By definition \eqref{definitionQ1},  
\begin{align*}
    \langle v\cdot \mathbf{u}_*',g\rangle_{\mathbf{u}_*}&=D\, \mathbf{u}_*'\cdot \int_{\mathbb{R}^d} v \, g\,G_{*}\,dv- D\int_{\mathbb{R}^d}(v\cdot \mathbf{u}_*') v  \,G_{*}\,dv \cdot \mathbf{v}_g.
\end{align*}
Assume $\mathbf{v}_g\neq 0$ and $\mathbf{u}_*'\neq0$. Since $\mathbf{v}_g \cdot \mathbf{u}_*'=0$, we can consider an orthonormal basis $\{e_1,e_2,..,e_d\}$ with $e_1=\mathbf{u}_*'/|\mathbf{u}_*'|$ and $e_2=\mathbf{v}_g/|\mathbf{v}_g|$. In these coordinates, by writing~$v=\sum_{i=1}^d v_ie_i$, we have $v\cdot \mathbf{u}_*'=v_1 |\mathbf{u}_*'|$ and $v\cdot \mathbf{v}_g=v_2|\mathbf{v}_g|$. Since the function $$v_1\mapsto v_1 \,v_2 \,G_{*} \left(\sum_{i=1}^dv_ie_i\right)=  v_1 \,v_2 \exp \left\{-\frac{1}{D}\left[\frac{\alpha }{4} \left(\sum_{i=1}^d|v_i|^2 \right)^2-\frac{\alpha}{2}\left( \sum_{i=1}^d |v_i^2|\right)+ \sum_{i=1}^d |v_i-\mathbf{u}_{f}\cdot e_i|^2\right] \right\}$$ is odd for all $(v_2,v_3,...v_d) \in \mathbb{R}^{d-1}$ because $\mathbf{u}_{f}\cdot e_i=0$ for $i\geq 2$, we obtain 
\begin{align*}
    |\mathbf{v}_g||\mathbf{u}_*'|\int_{\mathbb{R}^d}(v\cdot \mathbf{u}_*') v  \,G_{*}\,dv \cdot \mathbf{v}_g=\int \cdots \int  \, dv_2 \cdots dv_d
\left(\int_{\mathbb{R}}v_2 v_1  \,G_{*} \left(\sum_{i=1}^dv_ie_i\right)dv_1\right)=0,
\end{align*}
thanks to the change of variable $(v_1,v_2,...v_d)\mapsto \sum_{i=1}^dv_ie_i$.
Since $\mathbf{u}_*'\cdot \mathbf{v}_g=0$ we conclude 
$$\langle v\cdot \mathbf{u}_*',g\rangle_{\mathbf{u}_*}=0.$$
Hence, 

\begin{align}
\label{pop}
    \langle v\cdot \mathbf{u}_*'(1+g),g\rangle_{\mathbf{u}_*}=\langle v\cdot \mathbf{u}_*'g,g\rangle_{\mathbf{u}_*}=D\,\mathbf{u}_*' \cdot \int_{\mathbb{R}^d}v\,g^2G_{*}\,dv- D\mathbf{v}_g \cdot\int_{\mathbb{R}^d}(v\cdot \mathbf{u}_*')\,v\,g\,G_{*}\, dv.
\end{align}
Let us compute $\nabla_\mathbf{u}\langle g,g \rangle_\mathbf{u} $.
Using \eqref{definitionQ1} and the expression \eqref{nalblau}, we obtain:
\begin{align}
\label{bigequa2}
    \nabla_\mathbf{u}\langle g,g \rangle_\mathbf{u}&=D\,\int_{\mathbb{R}^d}g^2\,\nabla_\mathbf{u} G_\mathbf{u}\,dv- \nabla_\mathbf{u}\left| \notag \int_{\mathbb{R}^d}g\,v\,G_\mathbf{u}\,dv\right|^2\\&=\int_{\mathbb{R}^d}g^2(v-\mathbf{u})G_\mathbf{u}\,dv+2  \int_{\mathbb{R}^d}(v \otimes v)\, \mathbf{v}_g \,g\, G_\mathbf{u}\,dv-2\,D |\mathbf{v}_g|^2\mathbf{u},
\end{align}
where $v\otimes v$ denotes the matrix $(v_iv_j)_{1\leq i,j\leq d}$.\\
Finally,
since $\mathbf{u}_*'\cdot \mathbf{u}=0$, replacing \eqref{pop} and \eqref{bigequa2} in \eqref{bigequa},
\begin{align*}
    \frac{1}{2}\frac{d}{dt}\langle g,g\rangle_{\mathbf{u}_*}&=-Q_{2,\mathbf{u}_*}[g] + D^2 \mathbf{v}_g \cdot  \int_{\mathbb{R}^d}g(\nabla g- \mathbf{v}_g)G_{{*}}\,dv\\& - \left[\,\mathbf{u}_*' \cdot \int_{\mathbb{R}^d}v\,g^2G_{*}\,dv - \mathbf{v}_g \cdot\int_{\mathbb{R}^d}(v\cdot \mathbf{u}_*')\,v\,g\,G_{*}\,dv\right]  \\&+\frac{1}{2}  \mathbf{u}_*'\cdot \int_{\mathbb{R}^d}g^2\,v\,G_{*}\,dv +  \mathbf{u}_*'\cdot\int_{\mathbb{R}^d}(v\otimes v) \,\mathbf{v}_g\, g\, G_{*}\, dv.
\end{align*}
Using that $\mathbf{v}_g\cdot (v\cdot \mathbf{u}_*')v=\mathbf{u}_*'\cdot (v\otimes v)\mathbf{v}_g$, we obtain the result.
\end{proof}

\end{document}
